\FloatBarrier
\section{本章小结}
\label{sec:krylov-methods-summary}

本章从定常迭代每轮丢弃历史信息这一缺点出发，把初始残差及其矩阵乘积张成的Krylov子空间保留下来，在逐步扩大的仿射空间$\bm x_0+\mathcal K_k$中选取近似。选一个近似等价于选一个在原点取值$1$的残差多项式；各方法的区别在于按什么准则选这个多项式，以及怎样用性质良好的基而不是病态的幂向量完成计算。由于系数随当前向量变化，收敛性都由各自的选解准则推出，而不是由固定误差矩阵的谱半径判断。

方阵上的方法按矩阵结构分工。对称正定矩阵用CG，它在当前子空间中最小化误差的能量范数，步长与组合系数都可由残差长度算出，短递推在精确算术下保持全部正交与共轭关系，误差界由$\sqrt{\kappa_2(\bm A)}$控制。一般方阵用Arnoldi正交化得到Hessenberg关系，GMRES在其上最小化原残差，残差单调不增，但基的存储与正交化工作随步数增长；重启限制了存储，也失去了跨周期的最优性，可能停滞。对称不定矩阵用MINRES，Lanczos三项递推与带状QR使它在固定存储下仍最小化原残差，收敛界只依赖特征值。

非对称矩阵若要固定存储，就要放弃残差最小化。BiCG用原矩阵与转置矩阵维持双正交和双共轭，系数公式与CG同形，但残差可能波动，配对量为零时递推中断。BiCGSTAB用固定影子向量代替影子序列，方向的组合系数由残差多项式的首项系数比推出，再加一次局部最小残差步；每轮只用$\bm A$，局部改进只相对于中间残差成立，既不保证整轮单调，也不消除中断。

长方或秩亏的问题由Golub--Kahan双对角化处理：交替作用$\bm A$与$\bm A^\top$，在两个空间各得一组标准正交基，右侧基张成$\mathcal K_k(\bm A^\top\bm A,\bm A^\top\bm b)$，却从不形成$\bm A^\top\bm A$。LSQR在其中最小化原残差，精确算术下等价于正规方程上的CG；LSMR最小化正规残差，等价于正规方程上的MINRES，需要两层正交变换和两套方向。两者都取零初值，到达正规方程时选出最小范数解$\bm A^\dagger\bm b$；它们避免了形成Gram矩阵时的信息损失，但收敛快慢仍由非零奇异值平方的分布决定，原问题的病态性也不因算法而消失。

预条件用一个易解的$\bm M$改变迭代面对的谱：对角缩放、SSOR与不完全分解分别复用前几章的公式，PCG则是对称正定变换系统上的CG，换回原变量后只需每轮多解一次$\bm M\bm z=\bm r$。选用时要同时计入构造与应用的代价，并核对变换是否保持了方法的前提和原来的求解目标：PCG与MINRES需要对称正定的$\bm M$，左预条件会改变残差的度量，右预条件在秩亏时可能改变最小范数的选解。

最后，所有方法都以真实残差决定是否成功：方阵问题检查原残差，不相容最小二乘问题检查正规残差而允许原残差非零。递推得到的残差估计只是复核的提示；达到步数上限、停滞和中断都要与达到精度区分开。残差小到什么程度才说明解可信，取决于问题的条件数，这一问题以及在同一目标下怎样比较本章与前几章的方法，留给第~\ref{chap:solution-error-regularization}章。
